\section{导读：为什么激光雷达创业史属于科技产业课}

本节先说明这期为什么进入广义 AI/互联网方向。EP111 不是模型公司访谈，而是汽车产业链硬科技创业样本。过去十年，中国新能源汽车从无到有，整车品牌很显眼，但背后还有传感器、芯片、供应链、车规制造和主机厂关系这些隐形环节。禾赛的激光雷达故事，展示的是硬科技公司怎样从实验室技术走到车规量产。

这期的关键词不是“流量”，而是“降本、订单、定价、客户倒戈、主机厂、资本和国家产业机会”。它回答一个问题：当中国科技创新从互联网模式创新走向硬科技前沿创新，技术创业者需要补哪些能力？

\begin{figure}[H]
\centering
\includegraphics[width=0.96\textwidth]{hesai-timeline.png}
\caption{禾赛 11 年创业路线：从实验室技术、第一笔大单，到进入汽车大本营。自制概念图，依据视频章节和 00:03:40--03:02:16 对谈内容整理。}
\end{figure}

\begin{knowledgebox}{读图：硬科技创业不是单点发明}
激光雷达公司要同时解决技术、成本、制造、客户验证、融资、定价和产业窗口。任何一个环节掉链子，技术都很难变成产业位置。
\end{knowledgebox}

\begin{importantbox}{本期核心命题}
硬科技创业的机会常常来自产业链和国家能力的叠加：单个技术突破不够，必须与供应链、主机厂、资本周期、量产能力和产业窗口共振。
\end{importantbox}

\begin{knowledgebox}{硬科技创业 vs 互联网创业}
\begin{tabularx}{\textwidth}{p{0.20\textwidth}p{0.34\textwidth}X}
\toprule
维度 & 互联网创业 & 硬科技创业 \\
\midrule
核心资产 & 产品、流量、网络效应 & 技术、供应链、制造、客户认证。 \\
迭代速度 & 快速上线、快速 A/B & 样机、测试、认证、量产周期长。 \\
失败成本 & 可回滚、可改版 & 物料、库存、质量和客户风险高。 \\
融资用途 & 增长、研发、团队扩张 & 研发、设备、供应链、交付和库存。 \\
关键拐点 & 用户增长和留存 & 客户定点、量产、成本曲线和复购。 \\
\bottomrule
\end{tabularx}
\end{knowledgebox}

\subsection{本章小结}

EP111 是硬科技创业课。它把激光雷达从一个传感器，放进新能源汽车产业链和中国硬科技创业的大背景中理解。

